\documentclass[11pt]{article}
\usepackage[a4paper,margin=1in]{geometry}
\usepackage{amsmath,amssymb,amsfonts,bm}
\usepackage{hyperref}
\usepackage{enumitem}
\usepackage{mathtools}
\usepackage{amsthm}

\newtheorem{theorem}{Theorem}
\newtheorem{proposition}{Proposition}
\newtheorem{lemma}{Lemma}
\newtheorem{corollary}{Corollary}
\theoremstyle{remark}
\newtheorem{remark}{Remark}

\title{A Shared Mass--Energy Field with Visible and Dark Coupling Sectors:\\A First Lagrangian Draft}
\author{First Draft}
\date{\today}

\begin{document}
\maketitle

\begin{abstract}
Dark matter is usually inferred as gravitating matter that does not significantly emit, absorb, or scatter electromagnetic radiation. This draft formulates a speculative but conventional-field-theory-compatible idea: visible matter and dark matter are not separate because one has mass and the other does not; rather, they are different coupling sectors of a shared mass--energy-generating field. Visible excitations couple both to the shared mass field and to electromagnetism, while dark excitations couple to the mass field but have zero or suppressed electromagnetic charge. We present a minimal Lagrangian, derive the stress-energy tensor, identify symmetry requirements for dark-sector stability, and state the observational constraints. The model is intentionally modest: in first form it resembles hidden-sector dark matter, but its interpretive aim is to make ``visibility'' a coupling property rather than a mass property.
\end{abstract}

\section{Motivation}
The key intuition is:
\begin{quote}
Mass and visibility are not the same property. A field excitation may gravitate while failing to couple appreciably to electromagnetism.
\end{quote}
This statement is already compatible with standard dark-matter reasoning. The novel target of this draft is to package it into a shared mass--energy field framework in which visible and dark particles are different coupling modes of a deeper sector.

\section{Related models}
Hidden-sector dark matter models commonly introduce particles that interact weakly with the Standard Model and primarily reveal themselves gravitationally \cite{Cheung2010}. Self-interacting dark matter can arise from non-Abelian hidden sectors \cite{Boddy2014}, and dissipative hidden sectors may contain dark photons or additional forces \cite{Foot2014}. The present model is not yet a replacement for those frameworks. It is a minimal coupling-language reformulation that may later be made distinct by adding predictive structure.

\section{Fields and symmetries}
Let
\begin{itemize}[leftmargin=2em]
    \item $g_{\mu\nu}$ be the spacetime metric,
    \item $A_\mu$ be the electromagnetic gauge field,
    \item $\Phi$ be a shared mass--energy scalar field,
    \item $\psi$ be a visible-sector fermion,
    \item $\chi$ be a dark-sector fermion.
\end{itemize}
The model assumes that $\psi$ carries electromagnetic charge while $\chi$ does not:
\begin{equation}
    q_\psi\neq 0,
    \qquad
    q_\chi=0.
\end{equation}
A stabilizing dark symmetry may be imposed:
\begin{equation}
    \chi\rightarrow -\chi,
    \qquad
    \psi\rightarrow\psi,
    \qquad
    \Phi\rightarrow\Phi.
\end{equation}
This $\mathbb Z_2$ symmetry prevents dark matter decay into visible matter if no lighter odd states exist.

\section{Action}
The action is
\begin{equation}
    S=\int d^4x\sqrt{-g}\,\mathcal L,
\end{equation}
with
\begin{align}
    \mathcal L
    &= \frac{c^4}{16\pi G}R
    -\frac14 F_{\mu\nu}F^{\mu\nu}
    -\frac12 g^{\mu\nu}\partial_\mu\Phi\partial_\nu\Phi
    -V(\Phi)
    \nonumber\\
    &\quad
    +\bar\psi\left(i\hbar c\gamma^\mu D_\mu-m_{\psi,0}c^2\right)\psi
    +\bar\chi\left(i\hbar c\gamma^\mu \nabla_\mu-m_{\chi,0}c^2\right)\chi
    \nonumber\\
    &\quad
    -y_\psi\Phi\bar\psi\psi
    -y_\chi\Phi\bar\chi\chi
    -\lambda_{\psi\chi}(\bar\psi\psi)(\bar\chi\chi).
    \label{eq:main-lagrangian}
\end{align}
Here
\begin{equation}
    F_{\mu\nu}=\partial_\mu A_\nu-\partial_\nu A_\mu,
\end{equation}
while
\begin{equation}
    D_\mu=\nabla_\mu+i\frac{q_\psi}{\hbar c}A_\mu.
\end{equation}
The dark field $\chi$ has no electromagnetic covariant derivative because $q_\chi=0$.


In four spacetime dimensions the operator $(\bar\psi\psi)(\bar\chi\chi)$ has mass dimension six in natural units, so $\lambda_{\psi\chi}$ has mass dimension $-2$. It may therefore be viewed as an effective coefficient generated by integrating out heavier mediator physics.


\section{Effective masses}
If $\Phi$ develops a background value
\begin{equation}
    \langle\Phi\rangle=v_\Phi,
\end{equation}
then the effective masses become
\begin{equation}
    m_\psi^{\rm eff}=m_{\psi,0}+\frac{y_\psi v_\Phi}{c^2},
    \qquad
    m_\chi^{\rm eff}=m_{\chi,0}+\frac{y_\chi v_\Phi}{c^2}.
\end{equation}
Thus both visible and dark sectors obtain inertia through the shared field, but only the visible sector couples directly to $A_\mu$.

\section{Field equations}
Variation with respect to $\Phi$ gives
\begin{equation}
    \Box\Phi-\frac{dV}{d\Phi}
    = y_\psi\bar\psi\psi+y_\chi\bar\chi\chi.
\end{equation}
Variation with respect to $A_\mu$ gives Maxwell's equation
\begin{equation}
    \nabla_\mu F^{\mu\nu}=J^\nu_\psi,
\end{equation}
where
\begin{equation}
    J^\nu_\psi=q_\psi c\bar\psi\gamma^\nu\psi.
\end{equation}
There is no $J^\nu_\chi$ term unless a portal or millicharge is added.

Variation with respect to the metric gives
\begin{equation}
    G_{\mu\nu}=\frac{8\pi G}{c^4}T^{\rm total}_{\mu\nu},
\end{equation}
where
\begin{equation}
    T^{\rm total}_{\mu\nu}=T^{\Phi}_{\mu\nu}+T^{A}_{\mu\nu}+T^{\psi}_{\mu\nu}+T^{\chi}_{\mu\nu}+T^{\rm int}_{\mu\nu}.
\end{equation}
The scalar contribution is
\begin{equation}
    T^{\Phi}_{\mu\nu}
    =\partial_\mu\Phi\partial_\nu\Phi
    -g_{\mu\nu}\left[\frac12(\partial\Phi)^2+V(\Phi)\right].
\end{equation}


\subsection{Integrating out the shared scalar}
Write
\begin{equation}
 \Phi=v_\Phi+\phi
\end{equation}
and expand the potential near its local minimum,
\begin{equation}
 V(\Phi)=V(v_\Phi)+\frac12m_\phi^2\phi^2+O(\phi^3).
\end{equation}
In the regime where typical momentum transfer satisfies $|q^2|\ll m_\phi^2$, the scalar equation is algebraic to leading order:
\begin{equation}
 \phi\simeq-\frac{y_\psi\bar\psi\psi+y_\chi\bar\chi\chi}{m_\phi^2}.
 \label{eq:shared-heavy-scalar-solution}
\end{equation}
Substitution back into the Lagrangian gives
\begin{equation}
 \Delta\mathcal L_{\rm eff}
 \simeq \frac{1}{2m_\phi^2}
 \left(y_\psi\bar\psi\psi+y_\chi\bar\chi\chi\right)^2.
 \label{eq:shared-integrated-out}
\end{equation}
The visible--dark cross term is therefore
\begin{equation}
 \Delta\mathcal L_{\psi\chi}
 \simeq \frac{y_\psi y_\chi}{m_\phi^2}
 (\bar\psi\psi)(\bar\chi\chi).
 \label{eq:shared-derived-portal}
\end{equation}
Thus the four-fermion portal in Eq.~\eqref{eq:main-lagrangian} need not be independent: its natural magnitude is $|\lambda_{\psi\chi}|\sim |y_\psi y_\chi|/m_\phi^2$, with the sign determined by the convention used for the interaction term.

\begin{proposition}[Shared-mediator correlation]
In the heavy-mediator limit, both the visible--dark contact interaction and the environmental mass shifts are controlled by the same product of Yukawa couplings and the same mediator mass scale. Therefore suppressing the contact portal by increasing $m_\phi$ simultaneously suppresses long-wavelength environment-dependent mass correlations.
\end{proposition}


\section{Overlap-coupling matrix}
The coupling structure may be summarized by
\begin{equation}
    \mathcal C=
    \begin{pmatrix}
        y_\psi & q_\psi & g_\psi \\
        y_\chi & 0      & g_\chi
    \end{pmatrix},
\end{equation}
where the columns represent coupling to the mass field $\Phi$, electromagnetism $A_\mu$, and gravity $g_{\mu\nu}$. In ordinary units, gravity is universal at the stress-energy level, so $g_\psi$ and $g_\chi$ are not freely chosen charges. They indicate that both sectors contribute to $T_{\mu\nu}$.

The core condition for invisible gravitating matter is
\begin{equation}
    y_\chi\neq 0,
    \qquad
    q_\chi=0,
    \qquad
    T^\chi_{\mu\nu}\neq 0.
\end{equation}

\section{Dark-matter phenomenology}
The dark-sector relic abundance depends on production mechanism. If the sectors are thermally connected through the portal $\lambda_{\psi\chi}$, the relic density obeys a Boltzmann equation of the form
\begin{equation}
    \frac{dn_\chi}{dt}+3Hn_\chi
    =-\langle\sigma v\rangle
    \left(n_\chi^2-n_{\chi,{\rm eq}}^2\right).
\end{equation}
If the coupling is too weak for thermal equilibrium, freeze-in or gravitational production may be more appropriate.

Direct detection is suppressed when
\begin{equation}
    \lambda_{\psi\chi}\approx 0,
    \qquad
    q_\chi=0,
\end{equation}
while gravitational effects remain through $T^\chi_{\mu\nu}$.

\section{Possible new contribution}
To become more than a relabeled hidden-sector model, the proposal needs a distinctive prediction. One route is to make $\Phi$ spatially responsive to matter density:
\begin{equation}
    \Box\Phi-\frac{dV}{d\Phi}=y_\psi\bar\psi\psi+y_\chi\bar\chi\chi,
\end{equation}
so that dark and visible mass distributions are correlated by a shared mediator. This could modify halo profiles, self-interactions, or structure formation.

Another route is a density-dependent effective mass:
\begin{equation}
    m_\chi^{\rm eff}(x)=m_{\chi,0}+\frac{y_\chi\Phi(x)}{c^2},
\end{equation}
which would make dark matter respond to scalar-field environments.


\subsection{Static environmental response and fifth-force scale}
For a quadratic expansion of the scalar potential and a static nonrelativistic source, Eq.~(the scalar field equation) reduces schematically to
\begin{equation}
 (-\nabla^2+m_\phi^2)\,\delta\Phi
 =-\left(y_\psi n_\psi+y_\chi n_\chi\right),
 \label{eq:shared-yukawa-poisson}
\end{equation}
where the scalar densities have been replaced by number densities in the nonrelativistic limit. A uniform region much larger than $m_\phi^{-1}$ therefore has
\begin{equation}
 \delta\Phi
 \simeq-\frac{y_\psi n_\psi+y_\chi n_\chi}{m_\phi^2},
 \label{eq:shared-uniform-shift}
\end{equation}
and the two effective masses shift together according to
\begin{equation}
 \delta m_i^{\rm eff}=\frac{y_i\delta\Phi}{c^2}.
 \label{eq:shared-correlated-mass-shift}
\end{equation}
This is an explicit realization of the proposed correlated visible--dark response.

In natural units, exchange of $\phi$ between two static particles gives the Yukawa potential
\begin{equation}
 V_{ij}(r)=-\frac{y_i y_j}{4\pi r}e^{-m_\phi r}.
 \label{eq:shared-yukawa-potential}
\end{equation}
Relative to Newtonian gravity, the unscreened strength is
\begin{equation}
 \alpha_{ij}=\frac{y_i y_j}{4\pi Gm_i m_j}.
 \label{eq:shared-fifth-force-ratio}
\end{equation}
Equations~\eqref{eq:shared-uniform-shift} and \eqref{eq:shared-fifth-force-ratio} expose the model's central parameter tradeoff quantitatively: a light scalar increases the range of correlated mass response but also increases the range over which a visible-sector fifth force must be controlled.

\subsection{Gauge-invariant electroweak completion of the visible Yukawa term}
If $\psi$ denotes a Standard Model chiral fermion above electroweak symmetry breaking, the low-energy operator $\Phi\bar\psi\psi$ can arise from the gauge-invariant dimension-five interaction
\begin{equation}
 \mathcal L_{\rm EW}
 \supset-\frac{\kappa_\psi}{\Lambda}
 \Phi\,\bar L H\psi_R+\mathrm{h.c.}
 \label{eq:shared-ew-completion}
\end{equation}
for the appropriate electroweak multiplets. After $H$ acquires $\langle H\rangle=v_H/\sqrt2$, this becomes the low-energy coupling with
\begin{equation}
 y_\psi^{\rm eff}=\frac{\kappa_\psi v_H}{\sqrt2\Lambda}.
\end{equation}
Hence the simple Lagrangian can be interpreted either as a low-energy EFT or as the infrared limit of a gauge-invariant portal.


\section{Constraints}
The model must satisfy:
\begin{enumerate}[leftmargin=2em]
    \item no large fifth force in the visible sector;
    \item no excessive dark-radiation contribution to $N_{\rm eff}$;
    \item consistency with galaxy rotation curves and lensing;
    \item consistency with cosmic microwave background structure formation;
    \item direct-detection and collider limits on portal interactions;
    \item stability of the scalar potential $V(\Phi)$.
\end{enumerate}

\section{Conclusion}
The shared mass--energy field framework formalizes the idea that visible and dark matter may differ by coupling structure rather than by the mere possession of mass. In minimal form it is close to hidden-sector dark matter. Its value depends on whether the shared field $\Phi$ produces a new, constrained, and testable relation between visible matter, dark matter, and gravitational structure.

\begin{thebibliography}{9}
\bibitem{Cheung2010}
C. Cheung, G. Elor, L. J. Hall, and P. Kumar,
``Origins of Hidden Sector Dark Matter I,''
\texttt{arXiv:1010.0022}.

\bibitem{Boddy2014}
K. K. Boddy, J. L. Feng, M. Kaplinghat, and T. M. P. Tait,
``Self-Interacting Dark Matter from a Non-Abelian Hidden Sector,''
\texttt{arXiv:1402.3629}.

\bibitem{Foot2014}
R. Foot,
``Dissipative hidden sector dark matter,''
\texttt{arXiv:1409.7174}.

\bibitem{Adler2023}
S. L. Adler,
``Hidden Sector Dark Matter Realized as a Twin of the Standard Model,''
\texttt{arXiv:2308.08107}.
\end{thebibliography}

\end{document}
